% --- Archon physics formalization source begin ---
% archon:physics
% archon:chemistry
% archon:covers IChO2026Problems/problem_icho_2026_t6_a4.lean
% archon:source-report reports/icho_2026/problem_icho_2026_t6_a4.source.json
% archon:problem-id icho_2026_t6
% archon:part-id T6-A4

\chapter{Icho Chemistry problem icho\_2026\_t6\_a4}
\label{ch:IChO2026Problems_problem_icho_2026_t6_a4}

\paragraph{Problem source.}
\#\# Source
58th International Chemistry Olympiad, Tashkent, Uzbekistan, 2026.
Official-materials index: https://www.icho2026.uz/problems
English problem PDF: ../icho\_2026\_source/raw/theory\_problem.pdf
English solution/rubric PDF: ../icho\_2026\_source/raw/theory\_solution.pdf
Problem PDF page 53; printed local question page 2; solution starts on PDF page 51.
The page PNG is primary visual evidence for formulas, structures, tables, and figures.

\#\# T6. Carbon Nanorings
T6. Carbon Nanorings 7\% In 2019, C18, the first example of a new class of carbon allotropes called cyclo[n]carbons (C𝑛) was detected. Cyclocarbons are monocyclic, all-carbon compounds with high strain energy and low stability. They can be generated and analysed by bond-resolved Atomic Force Microscopy (AFM). Cyclocarbons have interesting aromatic properties based on their π systems. C13 has triplet (3C13, spin S = 1), and singlet (1C13, spin S = 0) forms. 6.1 Based on Hückel’s rule, fill in the table with the number of distinct aromatic (A) and anti-aromatic (AA) 𝜋-systems present in each of C18, C16, 3C13, and 1C13. Fill in all blanks. Later, cyclo[n]carbons C10, C12, C14, C20, and C26 were generated by a combination of dehalogenation and retro-Bergman (1) reactions on a surface. AFM images taken during the synthesis of C14 are shown below: 6.2 Draw structures B-D. Note they may contain one or more unpaired electrons. The structure of A is given as an example. 6.3 The voltages applied during the AFM-mediated synthesis of C𝑛may vary. Using the given C–X bond energies, choose the possible halogen(s) in the halogenated reagent for C18 synthesis via a retro-Bergman reaction with a voltage of 2.5 V acting on the electrons. Bond (C–X) Bond Dissociation Energy / kJ mol−1 C–F 467 C–Cl 346 C–Br 290 C–I 228 In 2025, the first relatively stable cyclo[48]carbon was synthesised. The molecule was stabilised by catenation (forming interlocking rings) with macrocycle E and was first characterised by electrospray mass spectrometry in positive mode.

\#\# Current subquestion 6.4
In the mass spectrum obtained, four intense peaks were observed m/z: 591, 783, 879, and 1174. Suggest the identity of these ions. Use integer atomic masses and assume no fragmentation happened. The ion corresponding to m/z = 591 is given as an example.

\#\# Dataset metadata
IChO 2026; T6; raw rubric points=12; kind=theory; formalization\_ready=true.

\paragraph{Shared problem context.}
T6. Carbon Nanorings 7\% In 2019, C18, the first example of a new class of carbon allotropes called cyclo[n]carbons (C𝑛) was detected. Cyclocarbons are monocyclic, all-carbon compounds with high strain energy and low stability. They can be generated and analysed by bond-resolved Atomic Force Microscopy (AFM). Cyclocarbons have interesting aromatic properties based on their π systems. C13 has triplet (3C13, spin S = 1), and singlet (1C13, spin S = 0) forms. 6.1 Based on Hückel’s rule, fill in the table with the number of distinct aromatic (A) and anti-aromatic (AA) 𝜋-systems present in each of C18, C16, 3C13, and 1C13. Fill in all blanks. Later, cyclo[n]carbons C10, C12, C14, C20, and C26 were generated by a combination of dehalogenation and retro-Bergman (1) reactions on a surface. AFM images taken during the synthesis of C14 are shown below: 6.2 Draw structures B-D. Note they may contain one or more unpaired electrons. The structure of A is given as an example. 6.3 The voltages applied during the AFM-mediated synthesis of C𝑛may vary. Using the given C–X bond energies, choose the possible halogen(s) in the halogenated reagent for C18 synthesis via a retro-Bergman reaction with a voltage of 2.5 V acting on the electrons. Bond (C–X) Bond Dissociation Energy / kJ mol−1 C–F 467 C–Cl 346 C–Br 290 C–I 228 In 2025, the first relatively stable cyclo[48]carbon was synthesised. The molecule was stabilised by catenation (forming interlocking rings) with macrocycle E and was first characterised by electrospray mass spectrometry in positive mode.

\paragraph{Current subquestion.}
In the mass spectrum obtained, four intense peaks were observed m/z: 591, 783, 879, and 1174. Suggest the identity of these ions. Use integer atomic masses and assume no fragmentation happened. The ion corresponding to m/z = 591 is given as an example.

\paragraph{Recorded answer/context.}
Molecular weight of E is 590 (C40H34N2O3). Molecular weight of C48 is 576. The first signal is given and corresponds to [E+H]+. The catenane with one molecule of E would yield signal m/z 1161 which does not match with given signals. Hence we should consider the probability of di- and tri-catenation as well as dication and trication formation (question clearly states positive mode): m/z +1H+ +2H+ +3H+ C48E 1167 584 389.7 C48E2 1757 879 586.3 C48E3 2347 1174 783 Therefore, 783 and 1174 peaks correspond to tri-catenane [C48E3 + 3H]3+ and [C48E3 + 2H]2+ ions. Peak at 879 corresponds to dicatenane dication[C48E2 + 2H]2+. m/z 591 783 879 1174 [E + H]+ [C48E3+3H]3+ [C48E2+2H]2+ [C48E3+2H]2+ 4 points for each correct ionic species. No partial credit as other answers will not have correct m/z ratios. 12.0 pt

\paragraph{Figure/image paths.}
\begin{itemize}
\item ../icho\_2026\_source/image/T6\_page-2.png
\item ../icho\_2026\_source/image/T6\_page-1.png
\end{itemize}

\paragraph{Formalization target.}
create a compiling Lean file with sorry bodies at `IChO2026Problems/problem\_icho\_2026\_t6\_a4.lean`.
The Lean declarations must preserve every mathematical object, quantifier, hypothesis, side condition, approximation/error guarantee, and requested conclusion from the source statement.
Use Mathlib/Physlib/Chemistry names found through LeanExplore where available. Prefer the configured benchmark Base library; introduce a local definition only when the source genuinely requires an object that the environment does not expose.

\begin{theorem}[Icho Chemistry formalization target]
\label{thm:physics:icho_2026_t6_a4:target}
The assigned autoformalize agent should translate this IChO chemistry problem into Lean declarations in the covered file, with theorem and lemma proof bodies written as `by sorry`.
\end{theorem}
\begin{proof}
This is an autoformalization task, not a proof task. Produce faithful statements that can later be proved without weakening the source contract.
\end{proof}
% --- Archon physics formalization source end ---
